\chapter{The Four Levels of Description}
\label{ch:four-levels}

\section{A Map as Four Documents}

The first thing to understand about the studio is that it does not hold a map. It
holds four descriptions of one, at four different grains, each owned by a
different tool and stored separately. The documentation states it flatly:

\begin{quote}
A map is not one document. It is four, each describing the whole map at a
different grain, each owned by a different tool and stored separately.
\end{quote}

\Cref{tab:levels} is that list. Read it top to bottom as increasing resolution:
the board knows what things are and not where they are exactly; the world knows
where every block is and nothing about what any of it is for.

\begin{table}[htbp]
\caption{The four levels. Each is a description of the whole map.}
\label{tab:levels}
\small
\begin{tabularx}{\textwidth}{L{2.1cm} Y L{3.1cm} L{1.9cm}}
\toprule
\textbf{Level} & \textbf{What it states} & \textbf{Document} & \textbf{Tool} \\
\midrule
The board &
  Rectangles on a coarse cell grid, what each one is for, and where the
  objectives sit &
  \id{PlanModel} \newline (\id{*.plan.json}) & Plan \\
The ground &
  The real geometry at block resolution: outlines, heights, relief, paint, props &
  \id{SketchLayout} & Sketch \\
The play &
  Teams, spawns, protections, build regions, objectives and how they are captured &
  \id{MapIntent} & Configure \\
The map &
  The voxel world and the \id{map.xml} a PGM server loads &
  \id{VoxelWorld} + \id{MapXml} & --- (built) \\
\bottomrule
\end{tabularx}
\end{table}

\subsection*{These are grains, not drafts}

The temptation is to read the four as stages of completeness, a rough sketch
refined into a finished thing. That reading is wrong and leads directly
to authoring mistakes, so the documentation heads it off:

\begin{quote}
A plan is not a rough draft of a layout: it states things a layout cannot, such as
that this rectangle is a wool room, or that these two pieces share a defence wall,
and cannot state things a layout can, like a curve or a one-block step. The same is
true upward: a layout knows where every block of ground is and has no idea what
any of it is for. That is why a finished map needs all four and why no one of them
is `the' map.
\end{quote}

This is worth dwelling on because it explains the shape of everything downstream.
A wool room is a concept that only exists at the board level. Once the board
compiles, there is no wool room in the geometry. There is a polygon, at a height,
with a theme on it. Conversely a coastline is a concept that only exists
at the ground level; no arrangement of rectangles can state one. An agent that
tries to express a curve in the plan is fighting the model, and an agent that
tries to state an objective in the layout is doing the same thing from the other
side.

\subsection*{The flow is one-way}

\begin{quote}
A plan compiles into a layout and an intent; a layout rasterizes into a world; an
intent projects into the map document; the document writes out as \id{map.xml}.
Nothing reads back up. There is no path from geometry to a plan, and none from a
finished \id{map.xml} to an intent. Both omissions are deliberate.
\end{quote}

Nothing infers upward. A finished map cannot be decompiled into the plan that
would have produced it, and the studio does not try; the documentation's position
is that hypothesising what a finished map's plan would have been is a person's job
rather than a tool's. The practical consequence for an agent is that corrections
are made \emph{upstream and re-run}, not patched downstream. If the ground is
wrong, the thing to change is the document that produced it.

\begin{figure}[htbp]
\centering
\begin{tikzpicture}[
  font=\sffamily\small,
  doc/.style={draw=rule, line width=0.9pt, fill=wash, rounded corners=2pt,
              minimum height=11mm, minimum width=25mm, align=center, inner sep=4pt},
  built/.style={doc, fill=accentsoft!45, draw=accent!40},
  every path/.style={draw=ink60, line width=0.8pt, -{Stealth[length=4pt]}}
]
\node[doc] (plan) {plan};
\node[doc, above right=4mm and 16mm of plan] (layout) {layout};
\node[doc, below right=4mm and 16mm of plan] (intent) {intent};
\node[built, right=16mm of layout] (world) {world};
\node[built, right=16mm of intent] (document) {document};
\node[built, below right=9mm and 15mm of world, minimum width=30mm] (out) {what a\\ server loads};

\draw (plan) -- node[above, sloped, font=\sffamily\scriptsize, color=ink60] {compile} (layout);
\draw (plan) -- node[below, sloped, font=\sffamily\scriptsize, color=ink60] {compile} (intent);
\draw (layout) -- node[above, font=\sffamily\scriptsize, color=ink60] {finish} (world);
\draw (intent) -- node[below, font=\sffamily\scriptsize, color=ink60] {project} (document);
\draw (world) -- (out);
\draw (document) -- (out);
\end{tikzpicture}
\caption{The five hand-offs. The two branches rejoin at the end rather than
staying independent: building the world \emph{resolves} parts of the intent, and
the \id{map.xml} is rendered from that resolved copy.}
\label{fig:handoffs}
\end{figure}

\section{The Five Hand-Offs}

\Cref{fig:handoffs} is the whole pipeline. Each arrow is a call.

\paragraph{Plan to layout and intent.} \route{POST /api/plan/compile} turns one
plan into both downstream documents, and it is pure: the same plan compiles to the
same pair on the server and in the browser. One thing happens here that surprises
people reading a compiled board for the first time. Abutting pieces of equal
height fuse into single polygons, so what arrives at the ground level is a board
rather than a grid of rectangles. A plan of eleven rectangles can compile to four
shapes.

\paragraph{Layout onto a map.} \route{PUT /api/map/\{slug\}/sketch/from-plan}
merges rather than replaces. Themes, room shells, dressing and author-corrected
heights are carried onto the fresh geometry. Relief is the exception, and the
reason is worth stating because it is the clearest example in the system of the
studio refusing to guess: relief is keyed by group id, group identity is derived
from the geometry, and a recompile that re-fuses the board produces different
groups. Hand-authored terrain would have nowhere correct to land. So that case
answers \textbf{409}, one \id{SK1} finding per group it would orphan, and
\id{?force=true} accepts the loss. It is the author's call, not the server's.

\paragraph{Intent onto a map.} \route{PUT /api/map/\{slug\}/intent/from-plan}
carries much less: the authors and contributors the stored intent already held,
and nothing else. Two slices deliberately do not ride across. \id{islandTeams} is
a positional derivation, and carrying it would relabel territory rather than
preserve an answer. \id{symmetry} is left behind because its presence is what
switches the orbit expander on. A rebuild therefore clears both, and the World and
Teams phases have to be walked again.

\paragraph{Layout to world.} \route{POST /api/map/\{slug\}/sketch/finish}
rasterizes the layout into world geometry and moves the map to the configure
stage. This is the only stage transition the studio performs at runtime.

\paragraph{Intent to document to \id{map.xml}.} \route{PUT /api/map/\{slug\}/intent}
stores the intent and projects it into the PGM document (teams, kits, regions,
filters, apply-rules and spawns) in one idempotent pass.
\route{GET /api/map/\{slug\}/xml} renders that document;
\route{GET /api/map/\{slug\}/export} gives the world.

\section{The Three-Document Load}

An agent reading the five hand-offs might reasonably conclude that authoring a map
means making five calls in order. It does not, and this is the single most
consequential piece of practical knowledge in the whole pipeline.

\route{POST /api/map/from-documents} takes a plan, a layout and an intent
together, stores a whole map from them, and answers the slug it landed under. The
compile, the rasterize and the projection all run inside it. A map already stored
under that slug is replaced, because the documents name one map, so loading them
twice is a reload rather than a second map.

The documentation is explicit that this is the authoring call and not merely the
import one, and equally explicit about what the staged path costs an agent that
walks it instead: six calls for one store, a fresh slug originated on every
correction, and the need to know that the intent's projection lands after the
metadata write. The staged path is what the browser does. A driver never needs it.

The layout and the intent are required; the plan is not. A board drawn on a grid
has no plan, so a layout emitter states none. Omitting either of the required two
answers \textbf{400} naming which.

This call is also the way \emph{back} in, and that matters more than it sounds.
No route in the studio reads a \id{map.xml} at all: \route{POST /map/import-folder}
refuses a folder carrying one outright, and \route{import-url} extracts only the
region files. So a map authored against one studio could reach another only as a
world, arriving without its plan, its drawing or its intent, and never able to be
re-planned. The three documents carry more than the world does.

\begin{ruling}{Request Body Conventions}
Twenty-five of the studio's write endpoints take a bare document and fourteen take
a wrapper, so an author guessing gets it right slightly more than half the time.
The rule that resolves it: if the endpoint is asking \emph{about a document}, post
the document unwrapped; if it is asking to \emph{file a document under a name},
post the record that names it. A wrong guess is silent, because an unknown
property is dropped rather than reported. Posting \id{\{"style": \{...\}\}}
where a bare style was wanted answers \textbf{200} with a preview of the defaults.
\end{ruling}

\section{The Library}

Beside the four sits one more thing that is not a level at all: the library of
materials, themes, house parts and room styles. It knows nothing about maps, having
neither slug nor stage, and the tools that build worlds reach into it to pick a
recipe.
What they take is a \emph{copy}, so a library edit can never rebuild a map that
already shipped. A theme changed today does not alter a world exported last week.

\section{Stage and Layer}

A map row carries a \textbf{stage} and, separately, the \textbf{layers} it holds.
A stage is where a map has got to. A layer is a document it carries, and it keeps
carrying it: a map built from a plan still has its plan after it has been
sketched, configured and exported.

A stage is a progress marker and not a lock. A configured map may be re-planned,
and no endpoint refuses on the basis of the stage a map stands at. What the stage
is for is saying which of the moves already open is the one the author was about
to make.

The studio answers both questions at runtime.
\route{GET /api/map/\{slug\}/state} gives the stage, the artifacts, and the moves
they allow, each with its route. A move is offered because the documents it reads
are stored, not because the stage is right. Its pair,
\route{GET /api/map/\{slug\}/findings}, answers what is wrong with the map right
now from every gate the stored documents can reach, and names the gates it did
\emph{not} ask along with the route that pays for them. An agent that has lost
track of where it is asks those two.

\section{Responsibilities Held at a Single Level}

A few things belong to exactly one level, and an agent looking for them at the
wrong one will not find them.

Terrain paint and dressing are the ground level's alone. The rasterizer is what
makes the world, so a finish authored a level above it would be authored before
the thing it finishes exists. The observer spawn is the play level's alone: a plan
puts it at the origin at a computed height, which is why an unattended spectator
stands at $0,0$. Water lanes are the board level's alone; they survive a
save at the play level and generate their region, but nothing at that level draws
them. Room shells come from the ground level's theme phase rather than from the
play level, which is the one most often looked for in the wrong place.
